\section{Analysis and Design for Composite Robustness}
\label{sec:composite}

A controlled-complexity task family, an efficient corrective proxy, then a parameter-design
application.

\subsection{Class-Imbalance--Controlled Tasks}
\label{subsec:class-imbalance}

\texttt{sweep\_class\_imbalance} and \texttt{sweep\_sparse\_class\_offset} parameterize a
controlled-complexity axis: instead of varying a task's spatial frequency, we vary the road's
\emph{share of GT pixels} (0.4--3.1\% in \texttt{sweep\_class\_imbalance}, with the road entirely
missing from the prediction) or its \emph{offset} (\texttt{sweep\_sparse\_class\_offset}, road
share fixed at ${\sim}0.4\%$). In both sweeps, \texttt{cbhm} (harsh, harmonic) sits at exactly
$0.000$ for every non-trivial perturbation level, while \texttt{cbhm\_soft} tracks the dominant
(building) class instead, staying between $0.758$ and $1.000$ across the same ranges. As a scalar
summary, define the \textbf{harsh-composite collapse threshold}: the smallest perturbation
magnitude at which \texttt{cbhm} first reaches exactly 0. Empirically this threshold is
effectively immediate --- \texttt{cbhm}'s harmonic mean has essentially zero tolerance once one
class is fully wrong, regardless of how small that class's share of the scene is.

\subsection{Weighted Reduction as an Efficient Corrective Proxy}
\label{subsec:weighted-reduction}

\begin{align}
  \mathrm{cbhm\_soft} = w_{\mathrm{linear}} \cdot \mathrm{clDice}_{\mathrm{mean,weighted}} +
                         w_{\mathrm{polygon}} \cdot \mathrm{BF}_{\mathrm{mean,weighted}}
\end{align}
where $w_{\mathrm{linear}} = N_{\mathrm{linear}} / (N_{\mathrm{linear}} + N_{\mathrm{polygon}})$
is each type's share of total GT pixels (0.5/0.5 if both totals are 0).

\begin{lstlisting}[language=Python, caption={\texttt{metrics/unified.py:88-102}}]
cl_mean_weighted = (float(np.dot(cldice_weights, cldice_scores) / cl_total)
                     if cl_total > 0 else 0.0)
bf_mean_weighted = (float(np.dot(bf_weights, bf_scores) / bf_total)
                     if bf_total > 0 else 0.0)
type_total = cl_total + bf_total
w_cl = float(cl_total / type_total) if type_total > 0 else 0.5
w_bf = float(bf_total / type_total) if type_total > 0 else 0.5
cbhm_soft = w_cl * cl_mean_weighted + w_bf * bf_mean_weighted
\end{lstlisting}

The key structural move: \texttt{cbhm\_soft} is a \textbf{reweighting of already-computed
per-class scores}, not a redesign of the harmonic-mean composite itself. It costs nothing beyond
the per-class scores DTAF1/CBHM already compute. \textbf{Stated carefully, not oversold}:
\texttt{cbhm\_soft} is a weighted \emph{arithmetic} mean, not geometric --- a geometric mean still
collapses to exactly 0 whenever one input is 0, same as harmonic (formalized in
Appendix~\ref{app:proofs}, Sections~\ref{app:a3}--\ref{app:a4}).

\textbf{Worked example --- the real failure that motivated this fix.} On real SpaceNet imagery
(\texttt{composite\_vs\_submetric\_report.ipynb}, tile \texttt{Khartoum\_img371}), a 12px road
offset gave $\mathrm{cbhm}=0.000$ despite $\mathrm{dtaf1}=0.934$ (per-class road F1 $=0.867$,
precision $=0.872$, recall $=0.863$) --- the road is only ${\sim}0.1\%$ of image pixels, and
clDice collapsed entirely on this sparse, offset road, dragging the harmonic mean to exactly zero
even though the dominant building class was essentially perfect. \texttt{cbhm\_soft} recovers
only partially here ($0.000 \to \mathbf{0.285}$), because the road actually has \emph{more} GT
pixels than the building despite its tiny share of image \emph{area} --- area-weighting only
helps when the failing class is also the pixel-count minority. Contrast
\texttt{Khartoum\_img333}, where the building genuinely dominates by pixel count: $\mathrm{cbhm}$
$0.48$--$0.52 \to \mathrm{cbhm\_soft}$ $\mathbf{0.85}$, a clean recovery. This nuance is treated as
a real limitation, not a clean win --- expanded in Section~\ref{sec:discussion}.

\subsection{Choosing Per-Class Tolerance Radii}
\label{subsec:tolerance-radii}

\texttt{DEFAULT\_TOLERANCES = \{"road": 10, "building": 2\}} (pixels) were validated, not guessed,
against the offset and erosion sweeps in Section~\ref{subsec:motivation-results}: the road offset
sweep shows DTAF1 holding at $1.000$ through exactly the 10px boundary and softening immediately
past it, confirming the tolerance behaves as intended rather than being either so tight it rejects
correct predictions or so loose it never penalizes anything. This is a cheap, already-available
sweep used as the design proxy instead of an expensive full retrain (or, here, a full real-data
re-annotation) for every candidate tolerance value --- though PLEM's tolerances are fixed
physically-motivated constants ($d = \text{physical\_metres} / \text{GSD\_metres\_per\_pixel}$)
rather than a learned/optimized parameter; Section~\ref{sec:conclusion} lists learned/adaptive
tolerance radii as future work. Section~\ref{subsec:tolerance-loss}'s \texttt{ToleranceBandLoss}
reuses these same tolerance values directly via a shared \texttt{class\_config}.
